\small
\renewcommand{\arraystretch}{1.2}

\begin{longtable}{|p{3.5cm}|p{10cm}|}

\caption{Filtros de destino aplicados nos modelos staging}
\label{tab:staging_filtros} \\

\hline
\textbf{Modelo} & \textbf{Filtro de Destino Brasil} \\ \hline
\endfirsthead

\multicolumn{2}{c}{\textit{Continuação da Tabela \ref{tab:staging_filtros}}} \\ \hline
\textbf{Modelo} & \textbf{Filtro de Destino Brasil} \\ \hline
\endhead

\hline
\multicolumn{2}{r}{\textit{Continua na próxima página}} \\
\endfoot

\hline
\endlastfoot

stg\_ikarus\_tours &
Busca pela palavra ``Brasil'' (em português/espanhol) e também por ``Brasilien'' (Brasil em alemão). Adicionalmente, identifica referências indiretas ao país por meio de termos como ``Amazon'' (Amazônia) e ``Iguazu'' (Cataratas do Iguaçu). \\ \hline

stg\_journey\_latin\_america &
Classifica registros como destino Brasil quando a URL contém padrões como \texttt{/brazil/}. Também realiza busca textual pela palavra ``Brazil'' nas descrições dos pacotes. \\ \hline

stg\_transalpino &
Aplica filtro por meio de uma CTE intermediária, mantendo apenas registros cujo campo \texttt{destino} ou a \texttt{url} contenham a palavra ``brasil'', sem diferenciação entre maiúsculas e minúsculas. \\ \hline

stg\_turismo\_costanera &
Remove registros com título vazio e mantém exclusivamente aqueles cujo campo \texttt{local} possui o valor ``Brasil'', dispensando análise textual complementar. \\

\end{longtable}